\input{../tex-inputs/latex-korrekturansicht-vorspann}

\section[Arthur Schnitzler an Berta Zuckerkandl, 8. 1. 1912]{L03945 Arthur Schnitzler an Berta Zuckerkandl, 8. 1. 1912}
\nopagebreak\mylabel{L03945v}
\rehead{ }\normalsize\beginnumbering\briefempfaengerindex{Zuckerkandl, Berta@\textsc{Zuckerkandl, Berta}!zzzSchnitzler, Arthur@\emph{von Arthur Schnitzler}!1912-01-081@{8. 1. 1912}|(be}
\toendnotes[C]{\smallbreak\pagebreak[2]}
\correspDesc{Versand  durch Arthur Schnitzler am 8. 1. 1912 in Wien
\newline{}Erhalt  durch Berta Zuckerkandl im Zeitraum [8. 1. 1912 –
            11. 1. 1912?] in Wien}\toendnotes[C]{\smallbreak}
\Standort{DLA, HS.1985.1.2282.}
\physDesc{Brief, Durchschlag, 1 Blatt, 2 Seiten, 989 Zeichen
\newline{}Schreibmaschine
\newline{}Handschrift: roter Buntstift, lateinische Kurrent (\noindent{}beschriftet: »\uline{\textsc{Zuckerkandl}}« und »\textcolor{pink}{\textsc{Fr}{[}an{]}\textsc{k}{[}reich{]}}\oindex{Frankreich@\textbf{Frankreich}|pw}«, sechs Unterstreichungen)}\toendnotes[C]{\smallbreak}
\pstart
           \raggedleft{}{\pb}8. 1. 1912.\pend
           
\pstart\center{}Verehrte gnädige Frau.\pend\vspace{0.5em}
\pstart
           Aus dem \label{K_L03945-1v}\edtext{beiliegenden Schreiben}{\lemma{\textnormal{\emph{beiliegenden Schreiben}}}\Cendnote{\textnormal{nicht
            überliefert}}}\label{K_L03945-1},
          dessen gelegentliche freundliche Rücksendung ich erbitte, ersehen Sie, dass Frau \textcolor{blue}{Lefreve}\pwindex{Lefèvre, A. @\textsc{Lefèvre, A.}, \emph{Übersetzerin}|pw}{}\ledrightnote{\textcolor{blue}{A. Lefèvre}} das »\textcolor{green}{Weite
            Land}\pwindex{Schnitzler, Arthur 15. 5. 1862 Wien – 21. 10. 1931 ebd.@\textsc{Schnitzler, Arthur} (15. 5. 1862 Wien – 21. 10. 1931 ebd.), \emph{Schriftsteller, Mediziner}!weite Land. Tragikomödie in fünf Akten@\strich\emph{Das weite Land. Tragikomödie in fünf Akten}|pw}{}\ledrightnote{\textcolor{green}{Das weite Land. Tragikomödie in fünf Akten}}« nicht übersetzen will. Indess hat mir \label{K_L03945-2v}\edtext{auch Herr \textcolor{blue}{Remon}\pwindex{Rémon, Maurice 27.\,11.\,1861 Paris – 20.\,6.\,1945 Mérignac@\textsc{Rémon, Maurice} (27.\,11.\,1861 Paris – 20.\,6.\,1945 Mérignac), \emph{Übersetzer}|pw}{}\ledrightnote{\textcolor{blue}{Maurice Rémon}}}{\lemma{\textnormal{\emph{auch Herr Remon}}}\Cendnote{\textnormal{\textcolor{blue}{Maurice
              Remon}\pwindex{Rémon, Maurice 27.\,11.\,1861 Paris – 20.\,6.\,1945 Mérignac@\textsc{Rémon, Maurice} (27.\,11.\,1861 Paris – 20.\,6.\,1945 Mérignac), \emph{Übersetzer}|pwk} an \textcolor{blue}{Arthur Schnitzler},
              20. 12. 1911, vgl. Karl Zieger: \emph{Arthur
                Schnitzler et la France 1894–1938. Enquête sur une réception},
              Villeneuve d’Ascq: \emph{Presses Universitaires du
                Septentrion} 2012, S. 191.
            }}}\label{K_L03945-2}, derjenige der sich auf seine
          persönlichen zu \textcolor{blue}{Guitry}\pwindex{Guitry, Lucien 13.\,12.\,1860 Paris – 1.\,6.\,1925 ebd.@\textsc{Guitry, Lucien} (13.\,12.\,1860 Paris – 1.\,6.\,1925 ebd.), \emph{Schriftsteller, Schauspieler}|pw}{}\ledrightnote{\textcolor{blue}{Lucien Guitry}} berief, seine Zweifel
          hinsichtlich der Chancen meines \textcolor{green}{Stückes}\pwindex{Schnitzler, Arthur 15. 5. 1862 Wien – 21. 10. 1931 ebd.@\textsc{Schnitzler, Arthur} (15. 5. 1862 Wien – 21. 10. 1931 ebd.), \emph{Schriftsteller, Mediziner}!weite Land. Tragikomödie in fünf Akten@\strich\emph{Das weite Land. Tragikomödie in fünf Akten}|pwv}{}\ledrightnote{{$\rightarrow$}\emph{\textcolor{green}{Das weite Land. Tragikomödie in fünf Akten}}} für \textcolor{pink}{Paris}\oindex{Paris@\textbf{Paris}, \emph{Hauptstadt}|pw}{}\ledrightnote{\textcolor{pink}{Paris}} ausgedrückt und so werde ich
          wohl die \textcolor{green}{Uebersetzung}\pwindex{Schnitzler, Arthur 15. 5. 1862 Wien – 21. 10. 1931 ebd.@\textsc{Schnitzler, Arthur} (15. 5. 1862 Wien – 21. 10. 1931 ebd.), \emph{Schriftsteller, Mediziner}!Le Pays Inconnu@\strich\emph{Le Pays Inconnu}|pwv}{}\ledrightnote{{$\rightarrow$}\emph{\textcolor{green}{Le Pays Inconnu}}} Herrn \textcolor{blue}{Morisse}\pwindex{Morisse, Paul 11.\,3.\,1866 Rouen – 28.\,9.\,1946 Paris@\textsc{Morisse, Paul} (11.\,3.\,1866 Rouen – 28.\,9.\,1946 Paris), \emph{Übersetzer}|pw}{}\ledrightnote{\textcolor{blue}{Paul Morisse}} übertragen, wenn der es nicht auch vorzieht
          abzurücken. Aus einer \label{K_L03945-3v}\edtext{\textcolor{green}{Notiz}\pwindex{P. C. @\textsc{P. C.}, \emph{Journalist/Journalistin}!Bei Le Bargy@\strich\emph{Bei Le Bargy}|pwv}{}\ledrightnote{{$\rightarrow$}\emph{\textcolor{green}{Bei Le Bargy}}} im \textcolor{green}{Neuen Wiener Journal}\pwindex{Neues Wiener Journal@\emph{Neues Wiener Journal}|pw}{}\ledrightnote{\textcolor{green}{Neues Wiener Journal}}}{\lemma{\textnormal{\emph{Notiz … Journal}}}\Cendnote{\textnormal{\textcolor{blue}{P. C.}\pwindex{P. C. @\textsc{P. C.}, \emph{Journalist/Journalistin}|pwk}: \emph{\textcolor{green}{Bei Le
                Bargy}\pwindex{P. C. @\textsc{P. C.}, \emph{Journalist/Journalistin}!Bei Le Bargy@\strich\emph{Bei Le Bargy}|pwk}}. In: \emph{\textcolor{green}{Neues Wiener Journal}\pwindex{Neues Wiener Journal@\emph{Neues Wiener Journal}|pwk}}, Jg. 20,
              Nr. 6539, 6. 1. 1912, S. 3–4. Darin wird geschildert, wie \textcolor{blue}{Le Bargy}\pwindex{Le Bargy, Charles Gustave 28.\,7.\,1858 La Chapelle – 5.\,2.\,1936 Nizza@\textsc{Le Bargy, Charles Gustave} (28.\,7.\,1858 La Chapelle – 5.\,2.\,1936 Nizza), \emph{Schauspieler}|pwk} im \emph{\textcolor{brown}{Burgtheater}\orgindex{Burgtheater@Burgtheater|pwk}} eine Aufführung von \emph{\textcolor{green}{Das weite
              Land}\pwindex{Schnitzler, Arthur 15. 5. 1862 Wien – 21. 10. 1931 ebd.@\textsc{Schnitzler, Arthur} (15. 5. 1862 Wien – 21. 10. 1931 ebd.), \emph{Schriftsteller, Mediziner}!weite Land. Tragikomödie in fünf Akten@\strich\emph{Das weite Land. Tragikomödie in fünf Akten}|pwk}} besuchte und genoss, ohne deutsch zu können.}}}\label{K_L03945-3} und was wohl
          massgebender ist auch von \label{K_L03945-4v}\edtext{\textcolor{blue}{\textcolor{blue}{privater Seite}\pwindex{Trebitsch, Siegfried 22.\,12.\,1868 Wien – 3.\,6.\,1956 Zürich@\textsc{Trebitsch, Siegfried} (22.\,12.\,1868 Wien – 3.\,6.\,1956 Zürich), \emph{Schriftsteller, Übersetzer}|pwv}{}\ledrightnote{{$\rightarrow$}\emph{\textcolor{blue}{Siegfried Trebitsch}}}}\pwindex{Trebitsch, Siegfried 22.\,12.\,1868 Wien – 3.\,6.\,1956 Zürich@\textsc{Trebitsch, Siegfried} (22.\,12.\,1868 Wien – 3.\,6.\,1956 Zürich), \emph{Schriftsteller, Übersetzer}|pwv}{}\ledrightnote{{$\rightarrow$}\emph{\textcolor{blue}{Siegfried Trebitsch}}} habe ich erfahren,
          dass \textcolor{blue}{Le Bargy}\pwindex{Le Bargy, Charles Gustave 28.\,7.\,1858 La Chapelle – 5.\,2.\,1936 Nizza@\textsc{Le Bargy, Charles Gustave} (28.\,7.\,1858 La Chapelle – 5.\,2.\,1936 Nizza), \emph{Schauspieler}|pw}{}\ledrightnote{\textcolor{blue}{Charles Gustave Le Bargy}}}{\lemma{\textnormal{\emph{privater … Bargy}}}\Cendnote{\textnormal{Durch \textcolor{blue}{Siegfried Trebitsch}\pwindex{Trebitsch, Siegfried 22.\,12.\,1868 Wien – 3.\,6.\,1956 Zürich@\textsc{Trebitsch, Siegfried} (22.\,12.\,1868 Wien – 3.\,6.\,1956 Zürich), \emph{Schriftsteller, Übersetzer}|pwk}, 
                  vgl. A. S.: \emph{Tagebuch}, 6. 1. 1912.}}}\label{K_L03945-4} sich sehr
          lebhaft für das \textcolor{green}{Stück}\pwindex{Schnitzler, Arthur 15. 5. 1862 Wien – 21. 10. 1931 ebd.@\textsc{Schnitzler, Arthur} (15. 5. 1862 Wien – 21. 10. 1931 ebd.), \emph{Schriftsteller, Mediziner}!weite Land. Tragikomödie in fünf Akten@\strich\emph{Das weite Land. Tragikomödie in fünf Akten}|pwv}{}\ledrightnote{{$\rightarrow$}\emph{\textcolor{green}{Das weite Land. Tragikomödie in fünf Akten}}}
            interessiert{[}.{]} Er hat es hier gesehen. Da er nicht deutsch versteht,
          ist diese Interesse keine besondere Bedeutung beizulegen; immerhin wäre zu überlegen, ob
          man mit ihm nicht in eine persönliche Verbindung treten könnte, wenn er im
            März hier gastieren wird.\pend
           
\pstart
           Bald hoffe ich Gelegenheit zu haben Sie {\pb}wiederzusehen
          und bin mit herzlichen Grüssen{\\[\baselineskip]} Ihr sehr ergebener\pend
           \leftskip=0em{}{\vspace{1\baselineskip}}
\pstart
           \noindent{}Frau Berta Zuckerkandl, \textcolor{pink}{Wien}\oindex{Wien@\textbf{Wien}, \emph{Verwaltungsgebiet}|pw}{}\ledrightnote{\textcolor{pink}{Wien}}.\pend
           \selectlanguage{ngerman}\endnumbering\briefempfaengerindex{Zuckerkandl, Berta@\textsc{Zuckerkandl, Berta}!zzzSchnitzler, Arthur@\emph{von Arthur Schnitzler}!1912-01-081@{8. 1. 1912}|)be}\mylabel{L03945h}
\begin{anhang}
\end{anhang}\newcommand{\dateiname}{L03945}\newcommand{\titel}{Arthur Schnitzler an Berta Zuckerkandl, 8. 1. 1912}\newcommand{\editorInnen}{Herausgegeben von Jahnke, SelmaMüller, Martin Anton}\input{../tex-inputs/latex-korrekturansicht-abspann}
